\subsection{全息地址测度、局部化扩张与 \texorpdfstring{$ZG$}{ZG} 传递维数的刚性后果}\label{subsec:conclusion-holographic-localized-zg-rigidity}
在定理 \ref{thm:xi-time-part62d-hologram-image-cantor-homeomorphism} 已把全息地址像识别为严格 cylinder 分解的 \(B\)-进 Cantor 系统、定理 \ref{thm:xi-time-part62dc-hologram-bernoulli-exact-dimensional} 已给出 Bernoulli 推前测度的精确维数公式、定理 \ref{thm:xi-localized-basic-extension-strictly-nonsplit} 已把局部化数系的离散/对偶短正合列锁定为非平凡扩张、而定理 \ref{thm:xi-time-part62dgc-zg-matrix-euler-no-scalar-product} 与定理 \ref{thm:killo-zg-dirichlet-matrix-euler} 已分别排除 \(ZG\) 的标量 Euler 因子化并给出其 \(2\times 2\) 传递实现之后，还可在结论层进一步抽出如下七条后果。最先出现的两条把全部事件流地址压缩到 \(T_2(E_{\min})\cong \ZZ_2^2\) 中的一条闭秩 \(1\) 细丝，并说明所有连通 Lie 相位对这条细丝都只能看到有限前缀；其余诸条则分别闭合于地址测度维数、Haar--interior 二分、局部化扩张不分裂、最终周期地址点的有理分类，以及 \(ZG\) 的最小 transfer 维数。

\begin{theorem}[Gödel 地址的二进 Tate 细丝]\label{thm:conclusion-godel-dyadic-tate-filament}
设 \(E\) 为有限事件集，取单射
\[
\mathrm{code}:E\hookrightarrow \{0,\dots,M\}\subset \ZZ,
\qquad
B=2^L>M,
\]
并固定 primitive 向量
\[
v\in T_2(E_{\min})\setminus 2T_2(E_{\min}).
\]
对任意事件流
\[
\mathbf e=(e_1,e_2,\dots)\in E^{\NN}
\]
定义全息地址
\[
X(\mathbf e):=\sum_{t\ge 1}\mathrm{code}(e_t)B^{t-1}v\in T_2(E_{\min}).
\]
则存在唯一的 \(2\)-进整数
\[
\alpha(\mathbf e):=\sum_{t\ge 1}\mathrm{code}(e_t)B^{t-1}\in \ZZ_2
\]
使得
\[
X(\mathbf e)=\alpha(\mathbf e)\,v.
\]
从而地址像满足
\[
\mathcal X_E:=X(E^{\NN})
\subseteq
T_v:=\ZZ_2v
\subset T_2(E_{\min})\cong \ZZ_2^2,
\]
并且 \(T_v\) 是一个闭的秩 \(1\) 子模。特别地，完整事件流并不充满整个 Tate 模块，而只落在其中一条余维 \(1\) 的 \(2\)-进细丝上；并且映射
\[
\mathbf e\longmapsto \alpha(\mathbf e)\longmapsto X(\mathbf e)
\]
为单射。
\end{theorem}
\begin{proof}
由于 \(|B|_2<1\) 且 \(\mathrm{code}(e_t)\) 有界，级数
\[
\sum_{t\ge 1}\mathrm{code}(e_t)B^{t-1}
\]
在 \(\ZZ_2\) 中收敛，并定义出唯一的
\[
\alpha(\mathbf e)\in \ZZ_2.
\]
把 \(v\) 提出即可得
\[
X(\mathbf e)=\alpha(\mathbf e)\,v.
\]
于是 \(\mathcal X_E\subseteq T_v:=\ZZ_2v\)。

另一方面，标量乘法给出 \(\ZZ_2\to T_v\)、\(a\mapsto av\) 的连续同构，故 \(T_v\cong \ZZ_2\) 是 \(T_2(E_{\min})\cong \ZZ_2^2\) 的闭秩 \(1\) 子模。

最后，由 \(B=2^L>M\) 的 \(B\)-进展开唯一性可知，不同事件流给出不同的 digit 序列 \((\mathrm{code}(e_t))_{t\ge 1}\)，故 \(\mathbf e\mapsto \alpha(\mathbf e)\) 为单射；再合成 \(a\mapsto av\) 即得 \(\mathbf e\mapsto X(\mathbf e)\) 的单射性。证毕。
\end{proof}

\begin{theorem}[连通 Lie 相位对 Gödel 细丝的湮没]\label{thm:conclusion-godel-filament-connected-lie-blindness}
选取复统一化
\[
E_{\min}(\CC)\cong \TT^2.
\]
则定理 \ref{thm:conclusion-elliptic-2adic-solenoid-exact-sequence} 给出短正合列
\[
0\longrightarrow T_2(E_{\min})\longrightarrow \Sigma_2(E_{\min})\xrightarrow{\ \pi_0\ }\TT^2\longrightarrow 0.
\]
由定理 \ref{thm:conclusion-godel-dyadic-tate-filament}，地址细丝 \(T_v=\ZZ_2v\) 完全落在核 \(T_2(E_{\min})=\ker \pi_0\) 中。因而复合映射
\[
E^{\NN}\xrightarrow{\ X\ }T_v\hookrightarrow \Sigma_2(E_{\min})\xrightarrow{\ \pi_0\ }\TT^2
\]
恒等于零。换言之，全部无限事件流在连通 torus 相位中完全不可见。

更一般地，设 \(G\) 为任意有限维连通实 Lie 群，且
\[
f:T_v\longrightarrow G
\]
为连续群同态。则 \(f(T_v)\) 必为有限 \(2\)-群；从而 \(\ker f\) 为开子群。于是存在 \(n_0\ge 1\)，使得任何两条具有相同前 \(n_0\) 个编码数字的事件流，都在 \(f\circ X\) 下具有相同像。特别地，任何尊重地址加法律的连通 Lie 通道，都只能记录有限前缀，而不能忠实承载无限 Gödel 地址。
\end{theorem}
\begin{proof}
第一部分仅是定理 \ref{thm:conclusion-elliptic-2adic-solenoid-exact-sequence} 与定理 \ref{thm:conclusion-godel-dyadic-tate-filament} 的直接合成：因为
\[
T_v\subseteq T_2(E_{\min})=\ker \pi_0,
\]
故 \(\pi_0\circ X\equiv 0\)。

对于第二部分，\(T_v\cong \ZZ_2\) 是紧、全不连通、pro-\(2\) 的交换群。其连续像 \(f(T_v)\) 仍是紧全不连通群；又因为 \(f(T_v)\) 是有限维 Lie 群 \(G\) 的闭子群，故 \(f(T_v)\) 本身也是一个零维紧 Lie 群，只能是有限群。连续像保持 pro-\(2\) 性，因此这个有限群必为 \(2\)-群。

有限像的核在紧群中总是开子群。由 \(T_v\cong \ZZ_2\) 的开子群分类，存在 \(N\ge 0\) 使
\[
\ker f=2^N T_v.
\]
若两条事件流的前 \(n_0\) 个编码数字相同，则其地址差满足
\[
X(\mathbf e)-X(\mathbf e')
\in
B^{n_0}T_v
=
2^{Ln_0}T_v.
\]
取 \(n_0\) 使 \(Ln_0\ge N\)，便有
\[
2^{Ln_0}T_v\subseteq 2^N T_v=\ker f,
\]
从而 \(f(X(\mathbf e))=f(X(\mathbf e'))\)。证毕。
\end{proof}

\begin{theorem}[全息地址的 Bernoulli 推前测度是 exact-dimensional，且维数完全闭式]\label{thm:conclusion-holographic-bernoulli-exact-dimensional}
沿用定理 \ref{thm:xi-time-part62d-hologram-image-cantor-homeomorphism} 的记号。设 \(E\) 为有限事件集，
\[
\mathrm{code}:E\hookrightarrow \{0,\dots,M\}\subset \ZZ,
\qquad
B=2^L>M,
\]
并对任意事件流
\[
\mathbf e=(e_1,e_2,\dots)\in E^{\NN}
\]
定义全息地址
\[
X(\mathbf e):=\sum_{t\ge 1}\mathrm{code}(e_t)B^{t-1}v,
\qquad
\mathcal X_E:=X(E^{\NN})\subset \ZZ_2v.
\]
对任意概率向量
\[
\mathbf p=(p_a)_{a\in E},
\qquad
H(\mathbf p):=-\sum_{a\in E}p_a\log p_a,
\]
记 \(\nu_{\mathbf p}:=\bigotimes_{n\ge 1}\sum_{a\in E}p_a\delta_a\) 为 \(E^{\NN}\) 上的 Bernoulli 乘积测度，并令
\[
\mu_{\mathbf p}:=X_\sharp \nu_{\mathbf p}.
\]
则 \(\mu_{\mathbf p}\) 是 exact-dimensional 测度，并且对 \(\mu_{\mathbf p}\)-几乎处处的 \(x\in \mathcal X_E\) 都有
\[
\dim_{\mathrm H}(\mu_{\mathbf p})
=
\dim_{\mathrm{loc}}(\mu_{\mathbf p},x)
=
\frac{H(\mathbf p)}{\log B}.
\]
特别地，若 \(\mathbf p\) 为 \(E\) 上的均匀分布，则
\[
\dim_{\mathrm H}(\mu_{\mathrm{unif}})
=
\frac{\log |E|}{\log B}.
\]
\end{theorem}
\begin{proof}
由定理 \ref{thm:xi-time-part62d-hologram-image-cantor-homeomorphism}，\(X:E^{\NN}\to \mathcal X_E\) 为拓扑同胚，并且任意长度 \(n\) 前缀
\[
[\mathbf a]:=\{\mathbf f\in E^{\NN}:\ \mathrm{pref}_n(\mathbf f)=\mathbf a\}
\qquad (\mathbf a\in E^n)
\]
在 \(\mathcal X_E\) 中对应的像为
\[
X([\mathbf a])=X_n(\mathbf a)+B^n\mathcal X_E.
\]
再由定理 \ref{thm:xi-time-part9g-holographic-prefix-isometry-on-line} 的柱--球对应，若 \(x=X(\mathbf e)\) 且 \(\mathbf a=(e_1,\dots,e_n)\)，则
\[
X([\mathbf a])=\mathrm B_{\mathcal X_E}(x,B^{-n}).
\]
因此
\[
\mu_{\mathbf p}\bigl(\mathrm B_{\mathcal X_E}(x,B^{-n})\bigr)
=
\nu_{\mathbf p}([\mathbf a])
=
\prod_{t=1}^{n}p_{e_t}.
\]
记
\[
E_+:=\{a\in E:\ p_a>0\}.
\]
对 \(\nu_{\mathbf p}\)-几乎处处的 \(\mathbf e\)，所有坐标 \(e_t\) 都落在 \(E_+\) 中；因此随机变量
\[
\xi(\mathbf e):=-\log p_{e_1}
\]
在 \(\nu_{\mathbf p}\) 下可积，并且
\[
\EE_{\nu_{\mathbf p}}[\xi]=H(\mathbf p).
\]
对 \(\xi\) 应用大数定律可得
\[
-\frac1n\log \mu_{\mathbf p}\bigl(\mathrm B_{\mathcal X_E}(x,B^{-n})\bigr)
=
\frac1n\sum_{t=1}^{n}(-\log p_{e_t})
\longrightarrow
H(\mathbf p).
\]
又因为
\[
\log(B^{-n})=-n\log B,
\]
遂得
\[
\lim_{n\to\infty}
\frac{\log \mu_{\mathbf p}\bigl(\mathrm B_{\mathcal X_E}(x,B^{-n})\bigr)}{\log(B^{-n})}
=
\frac{H(\mathbf p)}{\log B}.
\]
这表明 \(\mu_{\mathbf p}\) exact-dimensional，且其 Hausdorff 维数即为该常数。
均匀情形下 \(H(\mathbf p)=\log|E|\)，故得到最后一式。证毕。
\end{proof}

\begin{corollary}[全息地址像的 Haar--interior 二分律]\label{cor:conclusion-holographic-haar-interior-dichotomy}
沿用定理 \ref{thm:conclusion-holographic-bernoulli-exact-dimensional} 的记号，把 \(\mathcal X_E\) 视为 rank-\(1\) 地址线
\[
T_v:=\ZZ_2 v
\]
中的闭子集，并记 \(\mu_{T_v}\) 为 \(T_v\) 的归一化 Haar 概率测度。则有严格二分：
\begin{enumerate}
  \item 若 \(|E|<B\)，则
  \[
  \mu_{T_v}(\mathcal X_E)=0,
  \qquad
  \operatorname{int}_{T_v}(\mathcal X_E)=\varnothing;
  \]
  \item 若 \(|E|=B\)，则必有
  \[
  \mathrm{code}(E)=\{0,1,\dots,B-1\},
  \]
  并且
  \[
  \mathcal X_E=T_v=\ZZ_2v.
  \]
\end{enumerate}
\end{corollary}
\begin{proof}
先设 \(|E|<B\)。由定理 \ref{thm:xi-time-part62d-hologram-image-cantor-homeomorphism}，长度 \(n\) 的不同前缀恰给出 \(|E|^n\) 个两两不交的 cylinder
\[
X_n(\mathbf a)+B^n\mathcal X_E
\qquad (\mathbf a\in E^n),
\]
而它们都落在各自不同的 \(B^nT_v\)-余类中。由于
\[
T_v/B^nT_v\cong \ZZ/ B^n\ZZ,
\]
每个深度 \(n\) 的 Haar 球质量都等于 \(B^{-n}\)。故
\[
\mu_{T_v}(\mathcal X_E)
\le
|E|^n B^{-n}
=
\left(\frac{|E|}{B}\right)^n
\qquad (n\ge 1).
\]
令 \(n\to\infty\) 即得 \(\mu_{T_v}(\mathcal X_E)=0\)。

若 \(\mathcal X_E\) 在 \(T_v\) 中有内点，则存在某个深度 \(n\) 的 \(B\)-进球
\[
x+B^nT_v
\]
完全落在 \(\mathcal X_E\) 中。但该球在下一层会分裂为恰好 \(B\) 个深度 \(n+1\) 的子球，而 \(\mathcal X_E\) 在任一深度 \(n\) cylinder 内只保留 \(|E|\) 个下一层子 cylinder。由于 \(|E|<B\)，不可能包含整个父球，矛盾。故 \(\operatorname{int}_{T_v}(\mathcal X_E)=\varnothing\)。

再设 \(|E|=B\)。由于 \(\mathrm{code}(E)\subseteq\{0,\dots,M\}\) 且 \(B>M\)，必有 \(M=B-1\)，并且
\[
\mathrm{code}(E)=\{0,1,\dots,B-1\}.
\]
另一方面，\(B=2^L\) 意味着每个 \(z\in \ZZ_2\) 都有唯一的 \(B\)-进展开
\[
z=\sum_{t\ge 1}c_t B^{t-1},
\qquad
c_t\in\{0,\dots,B-1\}.
\]
于是对任意 \(x=zv\in T_v\)，都可唯一取
\[
e_t:=\mathrm{code}^{-1}(c_t)\in E
\]
使
\[
x=\sum_{t\ge 1}\mathrm{code}(e_t)B^{t-1}v=X(\mathbf e).
\]
故 \(T_v\subseteq \mathcal X_E\)。反向包含显然成立，因此 \(\mathcal X_E=T_v\)。证毕。
\end{proof}

\begin{theorem}[局部化数系扩张在离散层与对偶层都严格不分裂]\label{thm:conclusion-localized-discrete-dual-extension-nonsplit}
设 \(S\subset\mathbb P\) 为非空有限素数集，并记
\[
G_S:=\ZZ[S^{-1}].
\]
由定理 \ref{thm:killo-localization-circle-dimension-prime-ledger-splitting}，存在互为 Pontryagin 对偶的短正合列
\[
0\longrightarrow \ZZ\longrightarrow G_S\longrightarrow \bigoplus_{p\in S}C_{p^\infty}\longrightarrow 0,
\]
\[
0\longrightarrow \prod_{p\in S}\ZZ_p\longrightarrow \widehat{G_S}\longrightarrow \TT\longrightarrow 0.
\]
则这两条扩张都不分裂。等价地，
\[
G_S\not\cong \ZZ\oplus \bigoplus_{p\in S}C_{p^\infty},
\qquad
\widehat{G_S}\not\cong \TT\times \prod_{p\in S}\ZZ_p.
\]
\end{theorem}
\begin{proof}
若离散层短正合列分裂，则有
\[
G_S\cong \ZZ\oplus \bigoplus_{p\in S}C_{p^\infty}.
\]
右端含有非零挠元，而
\[
G_S=\ZZ[S^{-1}]\subset \QQ
\]
作为有理数加法群的子群是无挠群，矛盾。故离散侧不分裂。

若对偶层短正合列分裂，则
\[
\widehat{G_S}\cong \TT\times \prod_{p\in S}\ZZ_p.
\]
Pontryagin 对偶在局部紧交换群范畴中给出精确反变等价，因此对偶化后会得到离散侧短正合列分裂，这与上一步矛盾。故对偶侧亦不分裂。证毕。
\end{proof}

\begin{theorem}[最终周期事件流与 \texorpdfstring{$B$}{B}-进有理点的完全等价]\label{thm:conclusion-holographic-eventual-periodic-badic-rationals}
\leanverified{paper\_conclusion\_holographic\_eventual\_periodic\_badic\_rationals}
沿用定理 \ref{thm:conclusion-holographic-bernoulli-exact-dimensional} 的记号，并定义
\[
\QQ_B:=
\left\{
\frac{m}{B^\ell(B^r-1)}:\ m\in\ZZ,\ \ell\ge 0,\ r\ge 1
\right\}
\subset \QQ_2.
\]
则对任意 \(x\in \mathcal X_E\subset \ZZ_2v\)，下列两条等价：
\begin{enumerate}
  \item \(x=X(\mathbf e)\) 对某条最终周期事件流 \(\mathbf e\in E^{\NN}\)；
  \item \(x\in \QQ_B v\)。
\end{enumerate}
更具体地，若
\[
\mathbf e=(e_1,\dots,e_\ell,\overline{f_1,\dots,f_r}),
\]
并记
\[
a_t:=\mathrm{code}(e_t),
\qquad
b_j:=\mathrm{code}(f_j),
\]
则
\[
X(\mathbf e)
=
\left(
\sum_{t=1}^{\ell} a_t B^{t-1}
+
B^\ell\frac{\sum_{j=1}^{r} b_j B^{j-1}}{1-B^r}
\right)v
\in
\QQ_B v.
\]
\end{theorem}
\begin{proof}
若 \(\mathbf e\) 最终周期，写作
\[
\mathbf e=(e_1,\dots,e_\ell,\overline{f_1,\dots,f_r}),
\]
则其地址级数可分解为有限前缀与周期尾和：
\[
X(\mathbf e)
=
\sum_{t=1}^{\ell}a_t B^{t-1}v
+
B^\ell\sum_{q\ge 0}\sum_{j=1}^{r}b_j B^{j-1+qr}v.
\]
把后项视为几何级数即可得到
\[
X(\mathbf e)
=
\left(
\sum_{t=1}^{\ell} a_t B^{t-1}
+
B^\ell\frac{\sum_{j=1}^{r} b_j B^{j-1}}{1-B^r}
\right)v,
\]
故 \(X(\mathbf e)\in \QQ_B v\)。

反过来，设 \(x\in \QQ_B v\cap \mathcal X_E\)。则存在 \(q\in \QQ_B\) 与 \(\mathbf e\in E^{\NN}\) 使
\[
x=qv=X(\mathbf e).
\]
由 \(q\in \QQ_B\) 的定义，\(q\) 是一个 \(B\)-进有理点，因此按标准 \(B\)-进长除法，其展开必最终周期；也即存在唯一的 digit 序列
\[
(c_t)_{t\ge 1}\in \{0,\dots,B-1\}^{\NN}
\]
使
\[
q=\sum_{t\ge 1}c_t B^{t-1},
\]
且 \((c_t)\) 最终周期。
另一方面，由定理 \ref{thm:xi-time-part62d-hologram-image-cantor-homeomorphism} 的单射性，\(x\) 在 digits 属于 \(\{0,\dots,B-1\}\) 的 \(B\)-进展开是唯一的；而 \(\mathbf e\) 给出的 digits 正是 \(\mathrm{code}(e_t)\in\{0,\dots,M\}\subset\{0,\dots,B-1\}\)。故
\[
\mathrm{code}(e_t)=c_t
\qquad (t\ge 1),
\]
从而事件流 \(\mathbf e\) 也最终周期。证毕。
\end{proof}

\begin{theorem}[\texorpdfstring{$ZG$}{ZG} 的最小 transfer 维数恰等于 \texorpdfstring{$2$}{2}]\label{thm:conclusion-zg-minimal-transfer-dimension-two}
记
\[
ZG=
\left\{
\prod_{k\ge 1}p_k^{z_k}:\ z_k\in\{0,1\},\ z_kz_{k+1}=0,\ z_k=0\ \text{对充分大的 }k
\right\}.
\]
对 \(\Re(s)>1\)，其 Dirichlet 生成函数
\[
D_{ZG}(s):=\sum_{n\in ZG}\frac{1}{n^s}
\]
满足以下 sharp 结论：
\begin{enumerate}
  \item 存在 \(2\times 2\) transfer 表示
  \[
  D_{ZG}(s)
  =
  \lim_{N\to\infty}
  \mathbf 1^\top
  \Biggl(
  \prod_{i=1}^{N}
  \begin{pmatrix}
  1&1\\
  p_i^{-s}&0
  \end{pmatrix}
  \Biggr)
  \binom10;
  \]
  \item 不存在任何 \(1\times 1\) 标量局部 Euler 因子族
  \[
  D_{ZG}(s)=\prod_{p}F_p(p^{-s}),
  \qquad
  F_p(T)\in 1+T\CC[[T]],
  \]
  能在保持 Dirichlet 系数仍为 \(1_{ZG}(n)\) 的前提下表示同一 Dirichlet 级数。
\end{enumerate}
因此 \(ZG\) 的最小 transfer 维数恰为 \(2\)。
\end{theorem}
\begin{proof}
第一条正是定理 \ref{thm:killo-zg-dirichlet-matrix-euler} 的内容。

若第二条失败，即存在 \(1\times 1\) 的标量 transfer 表示，则其展开后的 Dirichlet 系数必按素数局部因子独立相乘，从而系数函数
\[
1_{ZG}(n)
\]
必须是乘法函数。现在取相邻素数 \(p_i,p_{i+1}\)。由定义，单点支撑
\[
(z_i,z_{i+1})=(1,0),\qquad (0,1)
\]
都满足最近邻约束，因此
\[
p_i\in ZG,
\qquad
p_{i+1}\in ZG.
\]
但
\[
p_ip_{i+1}\notin ZG,
\]
因为这对应于
\[
(z_i,z_{i+1})=(1,1),
\]
违反约束 \(z_kz_{k+1}=0\)。于是
\[
1_{ZG}(p_i)\,1_{ZG}(p_{i+1})=1,
\qquad
1_{ZG}(p_ip_{i+1})=0,
\]
与乘法性矛盾。故不存在标量 transfer 表示。

既然维数 \(1\) 不可能而维数 \(2\) 已实现，最小 transfer 维数恰为 \(2\)。证毕。
\end{proof}

\noindent
因此，H.160--H.162 所涉及的两条支线在结论层都达到刚性闭合。其一，rank-\(1\) 全息地址系统并不充满整个 \(T_2(E_{\min})\cong \ZZ_2^2\)，而是严格落在闭秩 \(1\) 子模 \(T_v=\ZZ_2v\) 上；相应地，任何连通 Lie 相位或连通 Lie 通道都只能看见这条细丝的有限前缀投影，而无法忠实承载其无限地址内容。其二，在这条同一地址线上，Bernoulli 事件源被压缩为熵除以 \(\log B\) 的 exact-dimensional 测度，最终周期点被完全识别为 \(B\)-进有理点；与此同时，局部化数系的 prime-register 核与圆周商形成真正的不分裂扩张，而 \(ZG\) 的最近邻禁词又把解析结构强制为最小 \(2\)-状态 transfer，不能退化为独立标量 Euler 因子的乘积。于是，地址细丝、Lie 失明、扩张类与传递维数四者便在同一组有限可核验证书中统一起来。

\endinput
